\section{A strict spectral surplus}\label{sec:counting}

Fix the standard orthogonal complex structure $J_0$ on $\mathbb R^{2s}$.
For $a,q>0$, let
$\lambda_{l,1}(a,q)\le\cdots\le\lambda_{l,N_l}(a,q)$
be the eigenvalues, counted with their real multiplicities, of the positive
angular operator in \eqref{eq:angular-operator} for $G(a,q,J_0)$.
Define the corresponding nonnegative exponents by
\begin{equation}\label{eq:dwell-exponents}
 d_{l,i}(a,q)\bigl(d_{l,i}(a,q)+m-1\bigr)=\lambda_{l,i}(a,q),
 \qquad d_{l,i}(a,q)\ge0.
\end{equation}
We will find more than the Euclidean number of directions whose exponents
have mean strictly below a common integer. The subsequent construction
will realize these means by cycling the selected directions.

The Hopf-charge decomposition underlying this calculation is described
by Tanno~\cite[Lemmas~3.1 and~4.1]{Tanno}; we prove the exact
multiplicities and the distribution needed for the count.

\begin{lemma}[Hopf-weight distribution]\label{lem:hopf-distribution}
For $m=2s-1$ and $0\le b\le l$, the eigenvalue $(2b-l)^2$ of
$-D_{J_0}^2$ on $V_l$ has the following contribution to its multiplicity:
\begin{align}
 h_{l,b}
 &=\binom{s+b-1}{s-1}\binom{s+l-b-1}{s-1}
   -\binom{s+b-2}{s-1}\binom{s+l-b-2}{s-1} \label{eq:hopf-multiplicity}\\
 &=\frac{l+s-1}{s-1}
   \binom{b+s-2}{s-2}\binom{l-b+s-2}{s-2}.\notag
\end{align}
Terms with a missing bidegree in the first line are zero.
With probabilities $h_{l,b}/N_l$, the quantities $(2b-l)^2/l^2$
converge in distribution, as $l\to\infty$, to
\[
 Y_m\sim\operatorname{Beta}\left(\frac12,\frac{m-1}{2}\right).
\]
\end{lemma}

\begin{proof}
Complexify $V_l$ and decompose homogeneous polynomials into bidegrees
$(b,l-b)$ in $z_1,\ldots,z_s$ and their conjugates.
The operator $D_{J_0}$ acts on this bidegree by $\mathrm i(2b-l)$.
The contraction $\sum_j\partial_{z_j}\partial_{\bar z_j}$ maps onto
the polynomials of bidegree $(b-1,l-b-1)$: in the positive monomial
inner product in which $z^\alpha\bar z^\beta$ has squared norm
$\alpha!\beta!$, its adjoint is multiplication by
$\sum_j z_j\bar z_j$, which is injective.
Taking the difference of dimensions gives the first line of
\eqref{eq:hopf-multiplicity}; cancellation gives the second.
The contributions indexed by \(b\) and \(l-b\) belong to the same
eigenvalue of \(-D_{J_0}^2\) and are added; the middle contribution is counted
once. Since this operator is real symmetric, the complex dimension of each
complexified eigenspace equals the real dimension of the corresponding
original real eigenspace. There is no additional
factor of two.

Summing the second line by the binomial convolution identity gives
\begin{equation}\label{eq:harmonic-multiplicity-asymptotic}
 N_l=\frac{l+s-1}{s-1}\binom{l+2s-3}{2s-3}
 =\binom{m+l}{m}-\binom{m+l-2}{m}
 \sim\frac{2}{(m-1)!}\,l^{m-1}.
\end{equation}
More precisely, $h_{l,b}/N_l$ is a beta-binomial probability.
If $X$ has density proportional to $[x(1-x)]^{s-2}$ on $(0,1)$ and,
conditionally on $X$, $K_l$ is binomial with parameters $l,X$, then
integration of the binomial probability against this density gives
$\mathbb P(K_l=b)=h_{l,b}/N_l$.
Since
\[
 \mathbb E\left[\left(\frac{K_l}{l}-X\right)^2\right]
 =\frac{\mathbb E[X(1-X)]}{l}\le\frac{1}{4l},
\]
we have $K_l/l\to X$ in probability in this representation.
The change of variable $y=(2x-1)^2$ gives density proportional to
$y^{-1/2}(1-y)^{s-2}$, proving the assertion.
\end{proof}

For $0<v<1$, define the upper-tail quantile $y_m(v)\in(0,1)$ by
\[
 \mathbb P\bigl(Y_m>y_m(v)\bigr)=v,
 \qquad
 \mathcal J_m:=\int_0^1 y_m(v)\log(1/v)\,dv.
\]
The threshold $2/(m-1)$ in the next lemma comes from the horizontal
factor in the link curvature condition~\eqref{eq:link-margin}.

\begin{lemma}[A weighted quantile inequality]\label{lem:weighted-quantile}
For all sufficiently large odd $m$,
\begin{equation}\label{eq:weighted-quantile-gap}
 (m-1)\mathcal J_m>2.
\end{equation}
\end{lemma}

\begin{proof}
Put $r=m-1$. The beta density, equivalently polar coordinates for
independent standard normal variables, gives the representation
\[
 rY_m\overset{\mathrm d}=\frac{Z^2}{Z^2/r+W_r/r},
\]
where $Z$ is standard normal and $W_r$ is an independent sum of $r$
squared standard normal variables. Since
$\mathbb E[W_r/r]=1$ and $\operatorname{Var}(W_r/r)=2/r$,
we obtain $rY_m\Rightarrow Z^2$.
The limiting distribution is continuous and strictly increasing on
$(0,\infty)$, so its upper quantile $y_\infty(v)$ satisfies
$r y_m(v)\to y_\infty(v)$ for every $0<v<1$.
Fatou's lemma therefore gives
\begin{equation}\label{eq:normal-quantile-limit}
 \liminf_{m\to\infty}(m-1)\mathcal J_m
 \ge \int_0^1 y_\infty(v)\log(1/v)\,dv
 =\mathbb E\left[R^2\log\frac{1}{Q(R)}\right],
\end{equation}
where $R=|Z|$ and $Q(x)=\mathbb P(R>x)$.
The last equality follows since $Q(R)$ is uniform on $(0,1)$.
This expectation is finite: integration of the half-normal density on
$[x,x+1]$ gives
$Q(x)\ge\sqrt{2/\pi}\,e^{-(x+1)^2/2}$.

For $x>0$, integration of the same density gives the classical Mills
upper bound (see \cite[Section~4, p.~113]{Lee1992})
\[
 Q(x)=\sqrt{\frac2\pi}\int_x^\infty e^{-t^2/2}\,dt
 \le\sqrt{\frac2\pi}\,\frac{e^{-x^2/2}}{x}.
\]
Consequently, writing $A=\mathbb E[R^2\log R]$, we have
\begin{equation}\label{eq:normal-quantile-lower-bound}
 \mathbb E\left[R^2\log\frac1{Q(R)}\right]
 \ge\frac32+A+\frac12\log\frac\pi2.
\end{equation}
Integration by parts against the half-normal density yields
$\mathbb E[R^4\log R]=3A+1$.
Under the probability measure $d\nu=R^4\,d\mathbb P/3$, Jensen's
inequality for the logarithm applied to $1/R$ now gives
\[
 A+\frac13=\mathbb E_\nu[\log R]
 \ge-\log\mathbb E_\nu[1/R]
 =\log\frac3{\mathbb E[R^3]}
 =\log\frac3{2\sqrt{2/\pi}}.
\]
Combining this with \eqref{eq:normal-quantile-lower-bound} proves
\begin{equation}\label{eq:strict-normal-constant}
 \mathbb E\left[R^2\log\frac1{Q(R)}\right]
 \ge\frac76+\log\frac{3\pi}{4}>2.
\end{equation}
Here the final strict inequality can be checked without decimal
approximations. Indeed, $\arctan x>x-x^3/3$ for $x>0$ gives
\[
 \frac\pi4=\arctan\frac12+\arctan\frac13
 >\frac{505}{648}>\frac79,
\]
so $3\pi/4>7/3$. Strict convexity of $t\mapsto1/t$ gives the
midpoint bounds
\[
 \log\frac73
 =\int_1^{5/3}\frac{dt}{t}+\int_{5/3}^{7/3}\frac{dt}{t}
 >\frac{2/3}{4/3}+\frac{2/3}{2}=\frac56.
\]
Equations~\eqref{eq:normal-quantile-limit} and
\eqref{eq:strict-normal-constant} imply \eqref{eq:weighted-quantile-gap}.
\end{proof}

Before selecting the link, fix a smooth nondecreasing cutoff \(\chi\),
zero near \(( -\infty,2]\) and one near \([4,\infty)\), and put
\[
 I_\chi=\int_2^\infty\chi(t)(1+t)^{-5/4}\,dt>0.
\]
The radial construction will require only the extra smallness condition
\begin{equation}\label{eq:cutoff-smallness}
 \frac{-\log a_*}{I_\chi}\,
 \sup_{t>2}\chi(t)(1+t)^{-5/4}\leq\frac14.
\end{equation}
It depends on the fixed cutoff and link scale, before any finite degree
range or control family has been chosen. The freedom in the next lemma
allows this condition to be imposed at the outset.

\begin{lemma}[Spectral surplus]\label{lem:spectral-surplus}
There is an odd $m=2s-1\ge7$ for which the following choices are possible.
The parameters $a_*\in(0,1)$ and $q_*>1$ can be taken arbitrarily close
to $1$, with \eqref{eq:link-margin} holding at $a=a_*$ for every
$q\in[1,q_*]$.
For these parameters there are integers $k\ge2$, $L\ge2$, and
$0\le M_l\le N_l$ for $2\le l\le L$ such that, on abbreviating
$d_{l,i}=d_{l,i}(a_*,q_*)$,
\begin{equation}\label{eq:selection}
 \overline d_l:=\frac1{M_l}\sum_{i=1}^{M_l}d_{l,i}<k
 \quad\text{whenever }M_l>0,
 \qquad
 \sum_{l=2}^L M_l>\sum_{l=0}^k N_l.
\end{equation}
\end{lemma}

\begin{proof}
\noindent\textbf{Choosing the link.}
Choose an odd $m\ge7$ satisfying Lemma~\ref{lem:weighted-quantile}, and
fix
\[
 \frac2{m-1}<c<\mathcal J_m.
\]
For a small positive parameter $\varepsilon$, set
\begin{equation}\label{eq:dwell-parameters}
 q_*=1+\varepsilon,
 \qquad \alpha_*^2=1-c\varepsilon,
 \qquad a_*^2=q_*^{1/m}\alpha_*^2.
\end{equation}
For $q\ge1$, the right-hand side of \eqref{eq:link-margin} is
\[
 F(q)=q^{1/m}\left(1-\frac{2(q-1)}{m-1}\right),
 \qquad
 F'(q)=\frac{m+1}{m(m-1)}q^{1/m-1}(1-2q)<0.
\]
At $q=q_*$, the required strict inequality follows from
$c>2/(m-1)$; it consequently holds throughout $[1,q_*]$.
In particular $a_*^2<F(q_*)<1$.

\medskip\noindent\textbf{Limiting means.}
By \eqref{eq:angular-operator}, the eigenvalues at the dwell parameters
are the numbers
\[
 \alpha_*^{-2}\left[B_l-
       \left(1-\frac1{q_*}\right)(2b-l)^2\right]
\]
with the multiplicities in Lemma~\ref{lem:hopf-distribution}.
Thus their increasing order corresponds to the decreasing order of the
Hopf weights $(2b-l)^2$.
If $d(d+m-1)=\lambda$ and $d\ge0$, then
$0\le\sqrt\lambda-d\le(m-1)/2$.
Since $B_l/l^2\to1$, Lemma~\ref{lem:hopf-distribution} and convergence
of quantiles to a continuous, strictly increasing distribution imply
\begin{equation}\label{eq:limiting-exponent-quantile}
 \frac{d_{l,\lceil vN_l\rceil}}{l}\longrightarrow
 f_\varepsilon(v):=(1-c\varepsilon)^{-1/2}
 \sqrt{1-\left(1-\frac1{q_*}\right)y_m(v)}
 \qquad(0<v<1).
\end{equation}
These rescaled exponents are uniformly bounded, so averaging the first
$\lceil uN_l\rceil$ terms gives, for every $0<u<1$,
\begin{equation}\label{eq:limiting-mean-exponents}
 \frac1{l\lceil uN_l\rceil}
       \sum_{i=1}^{\lceil uN_l\rceil}d_{l,i}
 \longrightarrow
 g_\varepsilon(u):=\frac1u\int_0^u f_\varepsilon(v)\,dv.
\end{equation}
One may see this directly by integrating the empirical quantile in
\eqref{eq:limiting-exponent-quantile}; the possible last fractional term
has size $O(N_l^{-1})$ after division by $l$.

Define
\[
 I(\varepsilon):=\int_0^1 g_\varepsilon(u)^{-m}\,du.
\]
All functions here are bounded away from zero for sufficiently small
$\varepsilon$. Uniformly for $0<v<1$,
\[
 f_\varepsilon(v)
 =1+\frac\varepsilon2\bigl(c-y_m(v)\bigr)+O(\varepsilon^2).
\]
The same uniform remainder remains valid after averaging over $(0,u)$.
Expansion of $g_\varepsilon(u)^{-m}$ and integration therefore give
\begin{align}
 I(\varepsilon)
 &=1+\frac{m\varepsilon}{2}
   \left[\int_0^1\frac1u\int_0^u y_m(v)\,dv\,du-c\right]
   +O(\varepsilon^2)\notag\\
 &=1+\frac{m\varepsilon}{2}(\mathcal J_m-c)
   +O(\varepsilon^2)>1
 \qquad\text{for all sufficiently small }\varepsilon>0.
 \label{eq:counting-first-variation}
\end{align}
The second equality uses the nonnegative integral identity
$\int_v^1du/u=\log(1/v)$.
Fix such an $\varepsilon$, as small as desired, from now on.

\medskip\noindent\textbf{Counting at integer thresholds.}
For each positive integer $k$ and $l\ge2$, let $M_l(k)$ be the largest
initial group of the ordered exponents whose mean is strictly less than
$k$, taking $M_l(k)=0$ if no group qualifies.
Only finitely many degrees contribute: the displayed eigenvalue formula
bounds all $\lambda_{l,i}$ below by a positive constant times $l^2$,
so $\min_i d_{l,i}$ grows at least linearly in $l$.
The initial-group means are nondecreasing, and hence
\begin{equation}\label{eq:counting-layer-cake}
 M_l(k)=N_l\int_0^1
 \mathbf 1\left\{
   \frac1{\lceil uN_l\rceil}
   \sum_{i=1}^{\lceil uN_l\rceil}d_{l,i}<k
 \right\}\,du.
\end{equation}
For fixed $u\in(0,1)$ and $\delta>0$,
\eqref{eq:limiting-mean-exponents} bounds the mean in this indicator by
$l(g_\varepsilon(u)+\delta/2)$ for every sufficiently large $l$.
It is therefore strictly below $k$ whenever
$l\le k/(g_\varepsilon(u)+\delta)$, including at the upper endpoint. Using
\eqref{eq:harmonic-multiplicity-asymptotic}, the sum of the corresponding
$N_l$, divided by $\sum_{l=0}^kN_l$, has limit
$(g_\varepsilon(u)+\delta)^{-m}$ as the integer $k\to\infty$.
Letting $\delta\downarrow0$ and applying Fatou's lemma to
\eqref{eq:counting-layer-cake} yields
\begin{equation}\label{eq:counting-liminf}
 \liminf_{\substack{k\to\infty\\ k\in\mathbb N}}
 \frac{\sum_{l=2}^\infty M_l(k)}{\sum_{l=0}^kN_l}
 \ge\int_0^1g_\varepsilon(u)^{-m}\,du
 =I(\varepsilon)>1.
\end{equation}
Discarding the finitely many lower degrees in this argument changes no
limit. Consequently every sufficiently large integer $k$ has a strict surplus.
Taking $L$ to be the largest degree with $M_l(k)>0$ and setting
$M_l=M_l(k)$ proves \eqref{eq:selection}.
\end{proof}

A selected prefix may be empty, a singleton, or all of \(V_l\), and may
cut a repeated real eigenspace. No gap at its last eigenvalue is required:
choose the indicated number of real orthonormal eigenvectors. The selected
span need not be invariant under any individual \(D_J\); the full fixed
space \(V_l\) remains invariant under every angular operator.

Finally, the sum appearing in \eqref{eq:selection} is exactly the
Euclidean comparison dimension in ambient dimension $n=m+1$:
\[
 \sum_{l=0}^kN_l
 =\binom{m+k}{k}+\binom{m+k-1}{k-1}
 =h_k(\mathbb R^{m+1}).
\]
To justify the Euclidean identification directly, let $u$ be a Euclidean
harmonic function of growth at most the integer $k$. A smooth radial kernel
$\psi_R(x)=R^{-(m+1)}\psi(x/R)$, supported in the ball of radius $R$ and
with integral $1$, reproduces $u$ by the mean-value property. For each
multi-index $\alpha$ of order $k+1$, moving derivatives onto the kernel
gives, at any fixed $x$ and for $R\ge1$,
\[
 |\partial^\alpha u(x)|
 \le \sup_{|y-x|\le R}|u(y)|\,
       \|\partial^\alpha\psi_R\|_{L^1}
 \le C_x R^kR^{-(k+1)}\longrightarrow0.
\]
Thus $u$ is a polynomial of degree at most $k$, and its homogeneous
components are harmonic. The converse follows from homogeneity.

For the rest of the construction, the spectral data mean an odd
\(m=2s-1\) with \(s\geq4\), a fixed cutoff \(\chi\), and
\(0<a_*<1<q_*\) satisfying \eqref{eq:link-margin} throughout
\([1,q_*]\) and \eqref{eq:cutoff-smallness}, together with finite
\(k,L\geq2\) and \(0\leq M_l\leq N_l\) satisfying
\eqref{eq:selection}. These are precisely the numerical hypotheses used
in Sections~\ref{sec:control}--\ref{sec:conclusion}; no later step uses
the asymptotic route by which the data were obtained.

The choices in Lemma~\ref{lem:spectral-surplus} are made in the order
$m,c,\varepsilon,k,L$. In particular any further requirement that $a_*$
be close to $1$ can be imposed when choosing $\varepsilon$, before the
integer cutoff is fixed.
